\chapter*{Resumen}
\addcontentsline{toc}{chapter}{Resumen}

El comercio electrónico (\emph{e-commerce}) argentino atraviesa un crecimiento sostenido, pero los emprendedores que operan tiendas en línea de pequeña escala carecen de herramientas analíticas que sean simultáneamente accesibles, integradas y prescriptivas. Las soluciones disponibles se limitan a reportes descriptivos, sin indicar qué acción tomar; las plataformas internacionales con capacidades prescriptivas presentan barreras de costo, idioma y adecuación al mercado local que las vuelven inaccesibles para este segmento.

El presente trabajo describe el diseño, la implementación y la evaluación de Decidata, un prototipo de sistema de recomendaciones accionables para tiendas en línea. Se integra en el panel de administración de Tiendanube y, mediante un adaptador multiplataforma, en WooCommerce. A partir de los datos históricos de ventas de cada tienda, el motor de análisis aplica suavizado exponencial simple y Holt-Winters para el pronóstico de demanda, segmentación de clientes mediante RFM y k-means, y un bandit multi-brazo con Thompson Sampling para priorizar las recomendaciones de mayor impacto esperado. Las recomendaciones se redactan en lenguaje natural con Gemini 3.5 Flash-Lite y se diferencian según el rubro comercial de la tienda: gastronomía, indumentaria o electrónica. % PENDIENTE-IMPL: B1
Además de recomendar, el sistema deja lista la acción, como un cupón o un precio promocional, para aplicarla con un clic y la confirmación del comerciante.

La investigación de usuarios se realizó mediante cinco entrevistas en profundidad con emprendedores de comercio electrónico y una sesión de brainstorming con perfiles técnicos, de producto y de marketing, entre ellos la dueña de un comercio electrónico. Con Design Thinking, las necesidades relevadas se tradujeron en historias de usuario y requerimientos formales.

El prototipo cuenta con pruebas automatizadas en el backend y en el frontend, y fue evaluado con datos generados por script para tres tiendas de demostración, una por rubro. El análisis estratégico y de viabilidad económico-financiera identifican una brecha de mercado e indicios de disposición a pagar en el segmento objetivo.

\bigskip
\noindent\textbf{Palabras clave:} comercio electrónico, recomendaciones prescriptivas, pronóstico de demanda, segmentación de clientes, bandit multi-brazo, Tiendanube, WooCommerce, inteligencia artificial.
